% --- Archon physics formalization source begin ---
% archon:physics
% archon:covers PhyXMiniProblems/problem_phyx_mini_0337.lean
% archon:source-report reports/phyx_mini/problem_phyx_mini_0337.source.json

\chapter{Physics problem phyx\_mini\_0337}
\label{ch:PhyXMiniProblems_problem_phyx_mini_0337}

\paragraph{Problem source.}
\#\# Physical scenario

The process \$abc\$ shown in the \$pV\$-diagram in \textbackslash{}textbf\{figure\} involves \$0.0175\$ mol of an ideal gas.

\#\# Question

How much work was done by or on the gas from \$a\$ to \$b\$? From \$b\$ to \$c\$? (c) If \$215\textasciitilde{}\textbackslash{}text\{J\}\$ of heat was put into the gas during \$abc\$, how many of those joules went into internal energy?

\#\# Answer choices

- A: A: 58J

- B: B: 53J

- C: C: 33J

- D: D: 63J

\#\# Figure caption (auxiliary; use the image as primary evidence)

The image is a graph plotting pressure (p) in atm on the y-axis and volume (V) in liters on the x-axis. It features three points labeled \textbackslash{}(a\textbackslash{}), \textbackslash{}(b\textbackslash{}), and \textbackslash{}(c\textbackslash{}). 

- Point \textbackslash{}(a\textbackslash{}) is at (2.0 L, 0.20 atm).
- Point \textbackslash{}(b\textbackslash{}) is at (2.0 L, 0.50 atm).
- Point \textbackslash{}(c\textbackslash{}) is at (6.0 L, 0.35 atm).

Arrows indicate transitions between these points:
- From \textbackslash{}(a\textbackslash{}) to \textbackslash{}(b\textbackslash{}) with a vertical arrow implying an increase in pressure.
- From \textbackslash{}(b\textbackslash{}) to \textbackslash{}(c\textbackslash{}) with a diagonal arrow indicating an increase in volume and a decrease in pressure.

The axes are labeled: \textbackslash{}(O\textbackslash{}) at the origin, \textbackslash{}(p\textbackslash{}) (atm) for pressure, and \textbackslash{}(V\textbackslash{}) (L) for volume. The grid serves as a background for better measurement precision.

\#\# Dataset metadata

Ideal Gases and Kinetic Theory; Physical Model Grounding Reasoning, Multi-Formula Reasoning

\paragraph{Recorded answer/context.}
B: B: 53J

\paragraph{Figure/image path.}
/root/proposal\_for\_physic/science-mango/phyx\_mini\_run/phyx\_data/test\_image/337.png

\paragraph{Formalization target.}
create a compiling Lean file with sorry bodies at `PhyXMiniProblems/problem\_phyx\_mini\_0337.lean`. The Lean declarations must preserve the physical quantities, dimensions or dimensional roles, figure labels, governing-law hypotheses, and final relation expressed by this problem.
Use Mathlib/Physlib names found through LeanExplore where available. If a physics API is missing, introduce faithful local abstractions rather than scalar placeholder aliases.

\begin{theorem}[Physics formalization target]
\label{thm:physics:phyx_mini_0337:target}
\lean{PhyXMiniProblems.ProblemPhyXMini0337.problem_phyx_mini_0337}
\uses{def:physics:phyx-mini-0337:phyxminiproblems-problemphyxmini0337-energyinjoules, def:physics:phyx-mini-0337:phyxminiproblems-problemphyxmini0337-idealgaspvprocess, def:physics:phyx-mini-0337:phyxminiproblems-problemphyxmini0337-workongasinjoules, def:physics:phyx-mini-0337:phyxminiproblems-problemphyxmini0337-hasstatedproblemdata, def:physics:phyx-mini-0337:phyxminiproblems-problemphyxmini0337-matchesprimarypvdiagram, def:physics:phyx-mini-0337:phyxminiproblems-problemphyxmini0337-hasphysicalpvparameters, def:physics:phyx-mini-0337:phyxminiproblems-problemphyxmini0337-satisfiesstraightpathboundaryworklaw, def:physics:phyx-mini-0337:phyxminiproblems-problemphyxmini0337-satisfiesfirstlawalongabc, def:physics:phyx-mini-0337:phyxminiproblems-problemphyxmini0337-displayedinternalenergyjoules, def:physics:phyx-mini-0337:phyxminiproblems-problemphyxmini0337-recordedanswerchoice, def:physics:phyx-mini-0337:phyxminiproblems-problemphyxmini0337-roundstonearestjoule}
The assigned autoformalize agent should translate this physics problem into Lean declarations in the covered file, with theorem and lemma proof bodies written as `by sorry`.
\end{theorem}
\begin{proof}
This is an autoformalization task, not a proof task. Produce statements that can later be proved by the physics prover without weakening the physical model.
\end{proof}

% --- Archon declaration topology begin ---
\section{Declaration topology}
\begin{definition}[Volume Quantity]
  \label{def:physics:phyx-mini-0337:phyxminiproblems-problemphyxmini0337-volumequantity}
  \lean{PhyXMiniProblems.ProblemPhyXMini0337.VolumeQuantity}
  A signed physical volume, with length-cubed dimension
\end{definition}
\begin{proof}
  This is the declared model definition of Volume Quantity.
\end{proof}
\begin{definition}[Pressure Quantity]
  \label{def:physics:phyx-mini-0337:phyxminiproblems-problemphyxmini0337-pressurequantity}
  \lean{PhyXMiniProblems.ProblemPhyXMini0337.PressureQuantity}
  A physical pressure using Physlib's pressure dimension
\end{definition}
\begin{proof}
  This is the declared model definition of Pressure Quantity.
\end{proof}
\begin{definition}[Energy Quantity]
  \label{def:physics:phyx-mini-0337:phyxminiproblems-problemphyxmini0337-energyquantity}
  \lean{PhyXMiniProblems.ProblemPhyXMini0337.EnergyQuantity}
  A signed physical energy using Physlib's energy dimension
\end{definition}
\begin{proof}
  This is the declared model definition of Energy Quantity.
\end{proof}
\begin{definition}[volume In Cubic Meters]
  \label{def:physics:phyx-mini-0337:phyxminiproblems-problemphyxmini0337-volumeincubicmeters}
  \lean{PhyXMiniProblems.ProblemPhyXMini0337.volumeInCubicMeters}
  \uses{def:physics:phyx-mini-0337:phyxminiproblems-problemphyxmini0337-volumequantity}
  Read a volume in SI cubic metres
\end{definition}
\begin{proof}
  This is the declared model definition of volume In Cubic Meters.
\end{proof}
\begin{definition}[volume In Liters]
  \label{def:physics:phyx-mini-0337:phyxminiproblems-problemphyxmini0337-volumeinliters}
  \lean{PhyXMiniProblems.ProblemPhyXMini0337.volumeInLiters}
  \uses{def:physics:phyx-mini-0337:phyxminiproblems-problemphyxmini0337-volumequantity, def:physics:phyx-mini-0337:phyxminiproblems-problemphyxmini0337-volumeincubicmeters}
  Read a volume in litres, using 1 m³ = 1000 L
\end{definition}
\begin{proof}
  This is the declared model definition of volume In Liters.
\end{proof}
\begin{definition}[pressure In Pascals]
  \label{def:physics:phyx-mini-0337:phyxminiproblems-problemphyxmini0337-pressureinpascals}
  \lean{PhyXMiniProblems.ProblemPhyXMini0337.pressureInPascals}
  \uses{def:physics:phyx-mini-0337:phyxminiproblems-problemphyxmini0337-pressurequantity}
  Read a pressure in pascals
\end{definition}
\begin{proof}
  This is the declared model definition of pressure In Pascals.
\end{proof}
\begin{definition}[pressure In Atmospheres]
  \label{def:physics:phyx-mini-0337:phyxminiproblems-problemphyxmini0337-pressureinatmospheres}
  \lean{PhyXMiniProblems.ProblemPhyXMini0337.pressureInAtmospheres}
  \uses{def:physics:phyx-mini-0337:phyxminiproblems-problemphyxmini0337-pressurequantity, def:physics:phyx-mini-0337:phyxminiproblems-problemphyxmini0337-pressureinpascals}
  Read a pressure in standard atmospheres
\end{definition}
\begin{proof}
  This is the declared model definition of pressure In Atmospheres.
\end{proof}
\begin{definition}[energy In Joules]
  \label{def:physics:phyx-mini-0337:phyxminiproblems-problemphyxmini0337-energyinjoules}
  \lean{PhyXMiniProblems.ProblemPhyXMini0337.energyInJoules}
  \uses{def:physics:phyx-mini-0337:phyxminiproblems-problemphyxmini0337-energyquantity}
  Read a signed physical energy in joules
\end{definition}
\begin{proof}
  This is the declared model definition of energy In Joules.
\end{proof}
\begin{definition}[joules Per Liter Atmosphere]
  \label{def:physics:phyx-mini-0337:phyxminiproblems-problemphyxmini0337-joulesperliteratmosphere}
  \lean{PhyXMiniProblems.ProblemPhyXMini0337.joulesPerLiterAtmosphere}
  The exact joule value of one litre-atmosphere
\end{definition}
\begin{proof}
  This is the declared model definition of joules Per Liter Atmosphere.
\end{proof}
\begin{definition}[Gas Model]
  \label{def:physics:phyx-mini-0337:phyxminiproblems-problemphyxmini0337-gasmodel}
  \lean{PhyXMiniProblems.ProblemPhyXMini0337.GasModel}
  The thermodynamic model named in the problem statement
\end{definition}
\begin{proof}
  This is the declared model definition of Gas Model.
\end{proof}
\begin{definition}[Process State]
  \label{def:physics:phyx-mini-0337:phyxminiproblems-problemphyxmini0337-processstate}
  \lean{PhyXMiniProblems.ProblemPhyXMini0337.ProcessState}
  The three labeled states shown on the pV diagram
\end{definition}
\begin{proof}
  This is the declared model definition of Process State.
\end{proof}
\begin{definition}[Process Leg]
  \label{def:physics:phyx-mini-0337:phyxminiproblems-problemphyxmini0337-processleg}
  \lean{PhyXMiniProblems.ProblemPhyXMini0337.ProcessLeg}
  The two directed legs of the process a → b → c
\end{definition}
\begin{proof}
  This is the declared model definition of Process Leg.
\end{proof}
\begin{definition}[Path Shape]
  \label{def:physics:phyx-mini-0337:phyxminiproblems-problemphyxmini0337-pathshape}
  \lean{PhyXMiniProblems.ProblemPhyXMini0337.PathShape}
  Geometric shape of a process leg in the pV plane
\end{definition}
\begin{proof}
  This is the declared model definition of Path Shape.
\end{proof}
\begin{definition}[Path Direction]
  \label{def:physics:phyx-mini-0337:phyxminiproblems-problemphyxmini0337-pathdirection}
  \lean{PhyXMiniProblems.ProblemPhyXMini0337.PathDirection}
  Directions of the two arrows as drawn in the primary figure
\end{definition}
\begin{proof}
  This is the declared model definition of Path Direction.
\end{proof}
\begin{definition}[Axis Label]
  \label{def:physics:phyx-mini-0337:phyxminiproblems-problemphyxmini0337-axislabel}
  \lean{PhyXMiniProblems.ProblemPhyXMini0337.AxisLabel}
  Literal labels printed on the two axes and at the origin
\end{definition}
\begin{proof}
  This is the declared model definition of Axis Label.
\end{proof}
\begin{definition}[Ideal Gas PVProcess]
  \label{def:physics:phyx-mini-0337:phyxminiproblems-problemphyxmini0337-idealgaspvprocess}
  \lean{PhyXMiniProblems.ProblemPhyXMini0337.IdealGasPVProcess}
  \uses{def:physics:phyx-mini-0337:phyxminiproblems-problemphyxmini0337-volumequantity, def:physics:phyx-mini-0337:phyxminiproblems-problemphyxmini0337-pressurequantity, def:physics:phyx-mini-0337:phyxminiproblems-problemphyxmini0337-energyquantity, def:physics:phyx-mini-0337:phyxminiproblems-problemphyxmini0337-gasmodel, def:physics:phyx-mini-0337:phyxminiproblems-problemphyxmini0337-processstate, def:physics:phyx-mini-0337:phyxminiproblems-problemphyxmini0337-processleg, def:physics:phyx-mini-0337:phyxminiproblems-problemphyxmini0337-pathshape, def:physics:phyx-mini-0337:phyxminiproblems-problemphyxmini0337-pathdirection, def:physics:phyx-mini-0337:phyxminiproblems-problemphyxmini0337-axislabel}
  The physical gas process and its figure transcription. workByGas is an independent dimensionful quantity for each leg. Neither it nor internalEnergyChangeAlongABC is defined to equal a requested answer
\end{definition}
\begin{proof}
  This is the declared model definition of Ideal Gas PVProcess.
\end{proof}
\begin{definition}[work On Gas In Joules]
  \label{def:physics:phyx-mini-0337:phyxminiproblems-problemphyxmini0337-workongasinjoules}
  \lean{PhyXMiniProblems.ProblemPhyXMini0337.workOnGasInJoules}
  \uses{def:physics:phyx-mini-0337:phyxminiproblems-problemphyxmini0337-energyinjoules, def:physics:phyx-mini-0337:phyxminiproblems-problemphyxmini0337-processleg, def:physics:phyx-mini-0337:phyxminiproblems-problemphyxmini0337-idealgaspvprocess}
  Work done on the gas, expressed using the opposite sign convention
\end{definition}
\begin{proof}
  This is the declared model definition of work On Gas In Joules.
\end{proof}
\begin{definition}[Has Stated Problem Data]
  \label{def:physics:phyx-mini-0337:phyxminiproblems-problemphyxmini0337-hasstatedproblemdata}
  \lean{PhyXMiniProblems.ProblemPhyXMini0337.HasStatedProblemData}
  \uses{def:physics:phyx-mini-0337:phyxminiproblems-problemphyxmini0337-energyinjoules, def:physics:phyx-mini-0337:phyxminiproblems-problemphyxmini0337-idealgaspvprocess}
  Problem-text data: 0.0175 mol of ideal gas and 215 J of heat supplied over the complete path. The unknown internal-energy change is not fixed here
\end{definition}
\begin{proof}
  This is the declared model definition of Has Stated Problem Data.
\end{proof}
\begin{definition}[Matches Primary PVDiagram]
  \label{def:physics:phyx-mini-0337:phyxminiproblems-problemphyxmini0337-matchesprimarypvdiagram}
  \lean{PhyXMiniProblems.ProblemPhyXMini0337.MatchesPrimaryPVDiagram}
  \uses{def:physics:phyx-mini-0337:phyxminiproblems-problemphyxmini0337-volumeinliters, def:physics:phyx-mini-0337:phyxminiproblems-problemphyxmini0337-pressureinatmospheres, def:physics:phyx-mini-0337:phyxminiproblems-problemphyxmini0337-idealgaspvprocess}
  Transcription of the primary raster. In particular, the image puts c on the 0.30 atm grid line, despite the auxiliary caption's 0.35 atm claim
\end{definition}
\begin{proof}
  This is the declared model definition of Matches Primary PVDiagram.
\end{proof}
\begin{definition}[Has Physical PVParameters]
  \label{def:physics:phyx-mini-0337:phyxminiproblems-problemphyxmini0337-hasphysicalpvparameters}
  \lean{PhyXMiniProblems.ProblemPhyXMini0337.HasPhysicalPVParameters}
  \uses{def:physics:phyx-mini-0337:phyxminiproblems-problemphyxmini0337-volumeinliters, def:physics:phyx-mini-0337:phyxminiproblems-problemphyxmini0337-pressureinatmospheres, def:physics:phyx-mini-0337:phyxminiproblems-problemphyxmini0337-idealgaspvprocess}
  Positivity conditions selecting physical pressures, volumes, and gas amount
\end{definition}
\begin{proof}
  This is the declared model definition of Has Physical PVParameters.
\end{proof}
\begin{definition}[Satisfies Straight Path Boundary Work Law]
  \label{def:physics:phyx-mini-0337:phyxminiproblems-problemphyxmini0337-satisfiesstraightpathboundaryworklaw}
  \lean{PhyXMiniProblems.ProblemPhyXMini0337.SatisfiesStraightPathBoundaryWorkLaw}
  \uses{def:physics:phyx-mini-0337:phyxminiproblems-problemphyxmini0337-volumeinliters, def:physics:phyx-mini-0337:phyxminiproblems-problemphyxmini0337-pressureinatmospheres, def:physics:phyx-mini-0337:phyxminiproblems-problemphyxmini0337-energyinjoules, def:physics:phyx-mini-0337:phyxminiproblems-problemphyxmini0337-joulesperliteratmosphere, def:physics:phyx-mini-0337:phyxminiproblems-problemphyxmini0337-idealgaspvprocess}
  Boundary-work law for a straight segment in a pV diagram. Since pressure varies linearly with volume on such a segment, W\_by = ((p\_start + p\_finish) / 2) * (V\_finish - V\_start). The last factor converts litre-atmospheres to joules. This is a general law for either leg and contains none of this problem's numerical answers
\end{definition}
\begin{proof}
  This is the declared model definition of Satisfies Straight Path Boundary Work Law.
\end{proof}
\begin{definition}[Satisfies First Law Along ABC]
  \label{def:physics:phyx-mini-0337:phyxminiproblems-problemphyxmini0337-satisfiesfirstlawalongabc}
  \lean{PhyXMiniProblems.ProblemPhyXMini0337.SatisfiesFirstLawAlongABC}
  \uses{def:physics:phyx-mini-0337:phyxminiproblems-problemphyxmini0337-energyinjoules, def:physics:phyx-mini-0337:phyxminiproblems-problemphyxmini0337-idealgaspvprocess}
  First law of thermodynamics for the complete route, with heat into the gas and work by the gas positive: ΔU = Q\_in - (W\_ab + W\_bc). The law relates three independent dimensionful-energy fields and does not state the requested numerical internal-energy change
\end{definition}
\begin{proof}
  This is the declared model definition of Satisfies First Law Along ABC.
\end{proof}
\begin{definition}[Answer Choice]
  \label{def:physics:phyx-mini-0337:phyxminiproblems-problemphyxmini0337-answerchoice}
  \lean{PhyXMiniProblems.ProblemPhyXMini0337.AnswerChoice}
  Labels of the four displayed internal-energy answer choices
\end{definition}
\begin{proof}
  This is the declared model definition of Answer Choice.
\end{proof}
\begin{definition}[displayed Internal Energy Joules]
  \label{def:physics:phyx-mini-0337:phyxminiproblems-problemphyxmini0337-displayedinternalenergyjoules}
  \lean{PhyXMiniProblems.ProblemPhyXMini0337.displayedInternalEnergyJoules}
  \uses{def:physics:phyx-mini-0337:phyxminiproblems-problemphyxmini0337-answerchoice}
  Joule readout printed beside each answer label
\end{definition}
\begin{proof}
  This is the declared model definition of displayed Internal Energy Joules.
\end{proof}
\begin{definition}[recorded Answer Choice]
  \label{def:physics:phyx-mini-0337:phyxminiproblems-problemphyxmini0337-recordedanswerchoice}
  \lean{PhyXMiniProblems.ProblemPhyXMini0337.recordedAnswerChoice}
  \uses{def:physics:phyx-mini-0337:phyxminiproblems-problemphyxmini0337-answerchoice}
  The answer label recorded in the supplied dataset metadata
\end{definition}
\begin{proof}
  This is the declared model definition of recorded Answer Choice.
\end{proof}
\begin{definition}[Rounds To Nearest Joule]
  \label{def:physics:phyx-mini-0337:phyxminiproblems-problemphyxmini0337-roundstonearestjoule}
  \lean{PhyXMiniProblems.ProblemPhyXMini0337.RoundsToNearestJoule}
  A real energy readout rounds to the displayed whole number of joules
\end{definition}
\begin{proof}
  This is the declared model definition of Rounds To Nearest Joule.
\end{proof}
\begin{lemma}[work Along Process Legs]
  \label{lem:physics:phyx-mini-0337:phyxminiproblems-problemphyxmini0337-workalongprocesslegs}
  \lean{PhyXMiniProblems.ProblemPhyXMini0337.workAlongProcessLegs}
  \uses{def:physics:phyx-mini-0337:phyxminiproblems-problemphyxmini0337-energyinjoules, def:physics:phyx-mini-0337:phyxminiproblems-problemphyxmini0337-idealgaspvprocess, def:physics:phyx-mini-0337:phyxminiproblems-problemphyxmini0337-workongasinjoules, def:physics:phyx-mini-0337:phyxminiproblems-problemphyxmini0337-matchesprimarypvdiagram, def:physics:phyx-mini-0337:phyxminiproblems-problemphyxmini0337-satisfiesstraightpathboundaryworklaw}
  The constant-volume leg does no work. On b → c, the average pressure is 0.40 atm and the volume increase is 4.0 L, so the gas does exactly 162.12 J = 4053/25 J of work (and the work on the gas is its negative)
\end{lemma}
\begin{proof}
  The conclusion follows from the governing relations and figure readouts named in the statement of work Along Process Legs.
\end{proof}
% --- Archon declaration topology end ---
% --- Archon physics formalization source end ---
